\section{Intégration du poste de traitement de surface}


\subsection{Présentation du poste de traitement de surface}

Le poste de traitement de surface est un système automatisé permettant de réaliser des opérations de traitement chimique ou galvanique sur des pièces. Il est composé d'un \textbf{pont roulant} motorisé qui déplace un \textbf{chariot} équipé d'un \textbf{crochet} pour immerger et transférer les pièces entre différentes \textbf{cuves de traitement}.

\subsubsection{Architecture matérielle}

\begin{itemize}
    \item \textbf{Automate WAGO 750-8012} : contrôleur PFC200 gérant le poste
    \begin{itemize}
        \item Adresse IP : \texttt{172.16.180.185}
        \item Communication : Ethernet/IP
        \item Type : Little Endian
    \end{itemize}

    \item \textbf{Chariot} (axe X - transversal) :
    \begin{itemize}
        \item Positionnement au-dessus des cuves
    \end{itemize}

    \item \textbf{Système de levage} (axe Z - vertical) :
    \begin{itemize}
        \item Crochet pour suspendre les pièces
        \item Montée/descente pour immersion dans les cuves
    \end{itemize}

    \item \textbf{Cuves de traitement} :
    \begin{itemize}
        \item Plusieurs cuves alignées contenant différents bains
        \item Positions fixes et connues
    \end{itemize}
\end{itemize}

% \subsubsection{Principe de fonctionnement}

% Le cycle de traitement se déroule selon les étapes suivantes :

% \begin{enumerate}
%     \item \textbf{Chargement} : prise de la pièce à la position de chargement
%     \item \textbf{Transfert} : déplacement du pont vers la première cuve
%     \item \textbf{Immersion} : descente du crochet pour immerger la pièce
%     \item \textbf{Temporisation} : maintien dans la cuve pendant le temps de traitement
%     \item \textbf{Extraction} : remontée du crochet
%     \item \textbf{Égouttage} : pause pour évacuation des liquides
%     \item \textbf{Transfert suivant} : répétition pour les cuves suivantes
%     \item \textbf{Déchargement} : dépose de la pièce traitée
% \end{enumerate}

\subsection{Rappel de la configuration actuelle}

\subsubsection{Premier coupleur : NOC\_ETHIP\_CONVOYEUR}

\begin{itemize}
    \item \textbf{Adresse IP} : \texttt{172.16.180.180}
    \item \textbf{Réseau} : \texttt{172.16.180.0/24} (classe B privée)
    \item \textbf{Capacité} : 128 mots en entrée, 128 mots en sortie
    \item \textbf{Équipement connecté} : PFC200 - Poste Guichet
    \begin{itemize}
        \item Adresse IP : \texttt{172.16.180.190}
        \item Données produites (T→O) : 12 octets 
        \item Données consommées (O→T) : 4 octets 
        \item Variables créées :
        \begin{itemize}
            \item \texttt{guichet\_IN} : type \texttt{T\_guichet\_IN}
            \item \texttt{guichet\_OUT} : type \texttt{T\_guichet\_OUT}
        \end{itemize}
    \end{itemize}
\end{itemize}


\subsubsection{Second coupleur : NOC\_ETHIP\_ROBOT}

\begin{itemize}
    \item \textbf{Adresse IP} : \texttt{192.168.0.1}
    \item \textbf{Réseau} : \texttt{192.168.0.0/24} (classe C privée)

    \item \textbf{Capacité} : 128 mots en entrée, 128 mots en sortie

    \item \textbf{Équipements connectés} (ordre d'ajout) :

    \begin{enumerate}
        \item \textbf{CS9 - Contrôleur robot Scara} (\texttt{192.168.0.10})
        \begin{itemize}
            \item Données produites (T→O) : 148 octets = 74 mots
            \item Données consommées (O→T) : 124 octets = 62 mots
            \item Variables créées :
            \begin{itemize}
                \item \texttt{cs9Scara\_IN} : type \texttt{T\_cs9Scara\_IN}
                \item \texttt{cs9Scara\_OUT} : type \texttt{T\_cs9Scara\_OUT}
            \end{itemize}
        \end{itemize}

        \item \textbf{M221 - Contrôleur E/S et HMI} (\texttt{192.168.0.11})
        \begin{itemize}
            \item Données produites (T→O) : 4 octets 
            \item Données consommées (O→T) : 40 octets 
            \item Variables créées :
            \begin{itemize}
                \item \texttt{m221Scara\_IN} : type \texttt{T\_m221Scara\_IN}
                \item \texttt{m221Scara\_OUT} : type \texttt{T\_m221Scara\_OUT}
            \end{itemize}
        \end{itemize}
    \end{enumerate}
\end{itemize}


\subsubsection{Services Ethernet/IP}

La communication avec le PFC200 du traitement de surface utilisera un service d'assemblage (IO-Scanning) similaire à celui mis en place pour le Guichet. Ce service échangera les données suivantes :

\begin{itemize}
    \item \textbf{Données produites par le pont} (T→O) : état du pont (positions, capteurs, mode en cours)
    \item \textbf{Données consommées par le pont} (O→T) : commandes vers le pont (consignes de mouvement, démarrage cycle)
\end{itemize}

\subsection{Questions}

\subsubsection{Configuration réseau}

\UPSTIquestion{Dessiner l'architecture réseau complète de la cellule robotisée incluant le poste de traitement de surface. Indiquer les adresses IP des différents équipements et automate.}

\subsubsection{Cartographie de la mémoire}
\UPSTIquestion{En analysant les adresses des appareils et en \textbf{justifiant votre réponse}, déterminer sur quel coupleur réseau (NOC) de l'automate M340 on configurera la liaison Ethernet/IP avec le contrôleur WAGO du traitement de surface ?}

Le module de traitement de surface produit et consomme les données suivantes, dans cet ordre :

\begin{center}
    \begin{tabular}{l|l}
        \textbf{Données   produites (T→O)}                              & \textbf{ Type }       \\\hline
        Etat du chariot \texttt{Arrêté, En mouvement, En Erreur, etc.}  & \textit{SINT}  \\
        Cuve de destination en cours                                    & \textit{USINT} \\
        Dernière position détectée                                      & \textit{USINT} \\
        Code Erreur Chariot                                             & \textit{SINT}  \\
        Etat de la pince (0. Entre deux, 1. En haut, 2. En bas)         & \textit{USINT} \\
        Code Erreur Pince                                               & \textit{SINT}  \\\hline\hline
        \textbf{Données consommées  (O→T)  }                            & \textbf{ Type}        \\\hline 
        Commande Chariot \texttt{Aller vers Cuve, Arrêter, etc.}        & \textit{SINT}  \\
        Cuve de destination                                             & \textit{USINT} \\
        Commande pince \texttt{Monter, Arrêter, Descendre}               & \textit{USINT} \\\hline
    \end{tabular}
\end{center}


\UPSTIquestion{
A partir des informations ci-dessus, établir la nouvelle cartographie mémoire. Elle devra contenir, au minimum : 
\begin{itemize}
    \item Les éléments de la configuration actuelle : 
    \begin{itemize}
        \item Les variables associées au premier coupleur (\texttt{172.16.180.180}) : adresses mémoires de base, taille mémoire en entrée/sortie, équipements connectés aves les adresses des variables associées
        \item Les variables associées au second coupleur (\texttt{192.168.0.1}) : adresses mémoires de base, taille mémoire en entrée/sortie, équipements connectés aves les adresses des variables associées
    \end{itemize}
    \item Les nouveaux éléments liés au poste de traitement de surface :
    \begin{itemize}
    \item Les adresses de début et de fin pour les entrées (T→O)
    \item Les adresses de début et de fin pour les sorties (O→T)
    \end{itemize}
\end{itemize}
}


\subsubsection{Déclaration des variables}
Deux types ont été créés dans le fichier donné : \texttt{TFromPontTS} et \texttt{TToPontTS}. Ils sont créé pour récupérer les données en provenance et à destination du module. 

\UPSTIquestion{
    Donner les déclaration des deux variables suivantes avec les adresses à leur affecter. 
\begin{itemize}
    \item \texttt{fromPontTS} de type \texttt{TFromPontTS} pour représenter les données produites
    \item \texttt{toPontTS} de type \texttt{TToPontTS} pour représenter les données consommées
\end{itemize}

Donner la syntaxe complète de déclaration avec les adresses mémoire appropriées.

\textit{Note : Les types \texttt{TFromPontTS} et \texttt{TToPontTS} sont fournis}
}

\subsubsection{Programmation de la communication}

\UPSTIquestion{
À quel moment du cycle automate la lecture des données produites par le module de traitement de surface doit-elle se faire ?
}

\UPSTIquestion{
À quel moment du cycle automate l'écriture des données consommées par le module doit-elle se faire ?
}


